% ============================================================================
% Appendix A6 — Carrying Capacity: Derivation and Flux-Threshold Reading
% ============================================================================
\section{Carrying Capacity: Derivation and Flux-Threshold Clarification}
\label{app:carrying}

\subsection{Derivation of Eq.~\eqref{eq:cmax}}

The depletion ratio \eqref{eq:gamma-ratio} compares the drain ceiling
$\delta O^{\gamma}$ (approached as $V/(V+\varepsilon)\to1$) with the
recovery ceiling $R$. Setting $\Gamma=1$ defines the critical complexity
\[
O_{\max} \;=\; \Bigl(\frac{R}{\delta}\Bigr)^{1/\gamma},
\]
and inverting the load kernel \eqref{eq:complexity} gives the capital level
at which the leader's own accumulation pushes complexity to that point:
\[
\Cmax
= \Biggl(\frac{(R/\delta)^{1/\gamma}-O_0}{\beta}\Biggr)^{1/\eta},
\]
which is Eq.~\eqref{eq:cmax}; if $(R/\delta)^{1/\gamma}\le O_0$ we set
$\Cmax=0$ (baseline load alone is unsustainable: the collapse-prone regime
of Definition~\ref{def:lowvitality}). At baseline,
$O_{\max}=\sqrt{150}=12.247$ and $\Cmax=(12.247-1)/0.25=44.99$. The
comparative statics \eqref{eq:cmax-statics} follow by inspection:
$\Cmax$ is decreasing in $\beta$ and $\delta$, increasing in $R$
(both pinned by \texttt{test\_carrying\_capacity\_monotone\_in\_beta} and
\texttt{\ldots\_in\_R}). The sensitivity to the load exponent is steep
(Table~\ref{tab:cmax-eta}).

\begin{table}[h]
\centering
\caption{Carrying capacity at baseline parameters as the load exponent
$\eta$ varies (Proposition~\ref{prop:carrying}).}
\label{tab:cmax-eta}
\vspace{0.4em}
\begin{tabular}{ccc}
\toprule
$\eta$ & $\Cmax$ & pinned by \\
\midrule
$1.0$ & $44.99$ & \texttt{test\_carrying\_capacity\_baseline\_equals\_45}\\
$1.5$ & $12.65$ & \texttt{test\_carrying\_capacity\_eta1\_5\_paper\_value}\\
$2.0$ & $6.71$  & \texttt{test\_carrying\_capacity\_eta2\_paper\_value}\\
\bottomrule
\end{tabular}
\end{table}

\subsection{What $\Cmax$ is: a flux threshold}

For $C>\Cmax$ we have $\Gamma>1$: the maximal drain exceeds the maximal
recovery, so $\dot V<0$ whenever $V/(V+\varepsilon)>R/(\delta O^{\gamma})$,
and the quasi-equilibrium $V_{\mathrm{qe}}(O)$ of
Proposition~\ref{prop:nullcline} falls into the depletion band
(asymptotically $V_{\mathrm{qe}}\approx R\varepsilon/(\delta O^{\gamma})$
in the strong-drain limit). $\Cmax$ therefore marks the capital level
beyond which no high-vitality operating point exists \emph{even
transiently at full recovery}---an energy-flux statement.

\subsection{What $\Cmax$ is not: a bifurcation}

Earlier drafts described $\Cmax$ as ``the critical threshold beyond which
sustainable equilibria cease to exist.'' That language is retracted, for
two independent reasons.

\begin{enumerate}
\item \emph{No fold anywhere.} By Lemma~\ref{lem:trace} and
Appendix~\ref{app:jacobian}, $\det J>0$ at every interior equilibrium for
all positive parameters. A saddle-node (fold) bifurcation requires
$\det J=0$; hence no interior equilibria are created or destroyed at
$\Cmax$ or anywhere else. The equilibrium count is governed solely by the
sign of $\Phi(O_0)$ (Theorem~\ref{thm:uniqueness}).
\item \emph{Trajectories cross it.} $\Cmax$ is not even a dynamical
barrier: the baseline trajectory from $(V_0,C_0)=(10,40)$ transiently
overshoots to $C\approx46.4>\Cmax=44.99$ before spiraling back
(Remark~\ref{rem:ctrap}). A fortiori it cannot serve as a trapping bound,
which is why the global-stability proof uses the distinct constant
$\Ctrap=102.67$ of Eq.~\eqref{eq:ctrap}
(Appendix~\ref{app:gas}).
\end{enumerate}

The corrected vocabulary is: $\Cmax$ = \emph{carrying capacity} (flux
threshold separating sustainable from collapse-prone \emph{regimes} as
parameters vary); $\Ctrap$ = \emph{trapping bound} (forward-invariance
ceiling used in Theorem~\ref{thm:gas}); and the two differ by more than a
factor of two at baseline.
